\
\documentclass[11pt]{article}
\usepackage[utf8]{inputenc}
\usepackage[T1]{fontenc}
\usepackage[margin=1in]{geometry}
\usepackage{hyperref}
\usepackage{enumitem}
\title{Block Coding with Bees}
\author{Piter Garcia}
\date{March 20, 2026}
\begin{document}
\maketitle
\section*{Quick Scan}
\begin{itemize}[leftmargin=*]
\item Grade: Grade 3
\item Packet focus: block coding, loops, and constraint-based planning in unplugged group tasks
\item Class blocks: Day 3 9:50-10:40 Lamanaco, Day 4 9:50-10:40 Lallucci, Day 5 9:50-10:40 Regelsberger, Day 5 2:15-3:05 Tandoi
\item Reference file: block-coding-with-bees.bib
\end{itemize}
\section*{School Planning Goals and Inclusion}
\begin{itemize}[leftmargin=*]
\item What is expected: make the CS learning target, proof, and revision path visible.
\item What we achieved: This packet is currently workbook-backed and should be enriched with classroom evidence as it becomes available.
\item What needs improvement: clearer proof, sharper assessment, and fewer hidden assumptions in directions.
\item How to improve: keep one artifact, one observation, and one revision note after reteaching.
\item Inclusion focus: visual modeling, chunked directions, role clarity, and flexible pacing supports.
\end{itemize}
\section*{Official References}
\begin{itemize}[leftmargin=*]
\item \url{https://csteachers.org/k12standards/interactive/}
\item \url{https://code.org/curriculum/elementary-school}
\end{itemize}
\end{document}
